\section{Finite traces of scalar--curvature forms}
\label{r11p:minimal-trace}

The trace will be constructed for the scalar--curvature module generated by products
\[
  u_0\,Du_1\cdots Du_k,\qquad D=d'd'',
\tag{MT1}\label{r11p:module-generators}
\]
where the $u_i$ are piecewise linear analytic potentials.

Throughout this section $K$ is a complete nontrivially valued field of
characteristic zero, and all analytic spaces are the analytifications of
normal integral separated $K$-varieties, with their ordinary topology.
All forms are Mihatsch actual delta-forms.  A PL scalar means a function
locally presented by a finite piecewise affine function of logarithms of
analytic units.  Its $D$ is a pure form of tridegree $(0,0,1)$.
Let $\mathscr M^k$ be the actual image sheaf of finite sums of
\eqref{r11p:module-generators}.  Relations in $\mathscr M^k$ are actual-form
relations; no chosen presentation is part of a section.

\subsection{Native degree under finite maps}

For a finite dominant map $q:Y\to X$ of generic degree $d$ and any compact
actual top form $\tau$ on $X$, these inputs give
\[
                            \int_Yq^*\tau=d\int_X\tau.
\tag{MT2}\label{r11p:native-degree}
\]
Indeed remove a proper closed subset outside which $q$ is finite flat of
constant degree $d$.  The support of $\tau$ avoids this subset by
Corollary 4.9; so does the support of its pullback above it.
Lemma 3.15 and the native skeletal integration formula apply on the
remaining open.    It also proves that finite dominant pullback on
actual forms is faithful: pair a putatively zero pullback with pulled-back
original-smooth complementary tests, use \eqref{r11p:native-degree}, and
apply Proposition 4.11.

\subsection{Descent of scalar PL functions through finite quotients}

\begin{lemma}[Ordinary local scalar PL descent]
\label{r11p:pl-descent}
Let $q:Y\to X$ be finite Galois between normal integral varieties, with
group $G$.  A $G$-invariant source PL scalar descends to a PL scalar on
$X$.  If finitely many invariant potentials on
$Y^{\rm an}\times(\mathbf G_m)^{\rm an}$ are jointly PL in fixed source
unit logarithms and the parameter $t=v(z)$, their descended potentials
are locally jointly PL in a fixed finite tuple of base-unit logarithms
and $t$.  The base tuple and base shrinking are independent of $t$.
\end{lemma}

\begin{proof}
Work around a whole finite analytic fiber.  Choose disjoint ordinary
neighborhoods of its points on which the finitely many source potentials
have moment presentations.  Properness permits shrinking the base so
that its full inverse image lies in their union.  The corresponding
source pieces are then clopen.  Analytic units can be pasted on these
pieces.  Add constant units with distinct valuations as markers; after
shrinking compact neighborhoods, their moment images lie in separated
compact polyhedral boxes.  Finite PL extension on these boxes
combines the presentations.  A smooth cutoff of a PL potential is
neither required nor used.

Include the full $G$-orbit of the resulting source-unit tuple.  Its
logarithms $\ell\in\mathbf R^s$ then carry a coordinate-permutation
$G$-action; add a constant unit if necessary so that $s\geq1$.
If an ambient PL presentation $\varphi(\ell)$ is
invariant only after restriction to the actual source moment image,
replace it by $|G|^{-1}\sum_{g\in G}\varphi(g\ell)$.  This is an
ambient invariant PL function with the same source restriction.
For ordered orbit potentials the sorting itself already gives this
ambient invariance.  The same operation leaves an independent
parameter fixed.  Put $d=|G|$ and $N=d(s-1)+1$.  For distinct integers $a_j$
take the Vandermonde directions
\[
                     m_j=(1,a_j,\ldots,a_j^{s-1})\in\mathbf Z^s.
\]
Define $J(\ell)$ by recording, for each $j$, the ordered multiset
$\{\langle m_j,g\ell\rangle:g\in G\}$.  It separates $G$-orbits.
In fact if $J(x)=J(y)$, for each $j$ one can choose $g_j$ with
$\langle m_j,x\rangle=\langle m_j,g_jy\rangle$.  One $g$ occurs for
at least $s$ indices.  The resulting polynomial of degree at most $s-1$
has $s$ roots, hence $x=gy$.

The map $J$ is proper: any $s$ independent Vandermonde directions bound
$\ell$ when all the displayed coordinates are bounded.  It is finite
piecewise linear.  On a common polyhedral subdivision its linearity
cells map affinely and injectively to their images; a noninjective
tangent direction would produce an infinite fiber.  Every invariant
PL scalar therefore factors as a PL scalar through $J$, after common
image subdivision and PL ambient extension.  The identical argument
applies to $J\times\operatorname{id}_{\mathbf R}$ and jointly PL
potentials with the fixed parameter $t$.

It remains to realize $J$ downstairs without an illicit nonlinear
pullback of forms.  For each source monomial $F_j=f^{m_j}$ take its monic
orbit polynomial
\[
                          \prod_{g\in G}(T-gF_j).
\]
Its coefficients are invariant analytic functions, hence analytic
functions on the normal finite quotient.  Algebraically the quotient
is $X$; analytically this follows on affinoids from finite coherent
base change and the Reynolds idempotent in characteristic zero.
With the convention that the coefficient of $T^i$ gives the point
$(i,v(a_i))$, the ordered root valuations are the negatives of the
lower Newton polygon's slopes, with multiplicity.  They are therefore
PL in the valuations of the coefficients that contribute to that
polygon.

At a fixed base point retain the nonzero coefficients and shrink so
that they are units.  The constant coefficient $a_0$ is a unit because
it is a product of source units; the leading coefficient is one.
Choose a uniform upper bound $B$ for $v(a_0)$ on the shrunken base.
Every coefficient vanishing at the fixed point can, by continuity and
a further common shrinking, be made to have valuation greater than
$\max(B,0)+1$.  Its point is then above the endpoint chord of the
Newton polygon and cannot affect the lower hull.  Such coefficients
can be omitted uniformly.
Thus $J$ is PL in finitely many base-unit logarithms.  These
coefficients involve source units alone, so their choice and shrinking
are independent of $t$.  Substitution is only a substitution into a
scalar PL function.  Taking its actual derivatives happens afterwards;
no nonlinear PL pullback of a general delta-form has been postulated.
\end{proof}

\subsection{Sorting gives a pure-curvature trace}

\begin{lemma}[Max/min identities]
\label{r11p:sorting}
For source PL scalars $u,v$, put $M=\max(u,v)$ and $m=\min(u,v)$.
For every $k\geq0$ there are actual identities
\[
 (DM)^k+(Dm)^k=(Du)^k+(Dv)^k,
\tag{MT3}\label{r11p:sorting-powers}
\]
\[
 M(DM)^k+m(Dm)^k=u(Du)^k+v(Dv)^k.
\tag{MT4}\label{r11p:sorting-energies}
\]
\end{lemma}

\begin{proof}
Write $h=u-v$, $h_+=\max(h,0)$, $h_-=\min(h,0)$, and
$A=Dh_+$, $B=Dh_-$.  Coefficientwise multiplication on support
polyhedra gives $h_+B=h_-A=0$: on one open side the appropriate
function is zero, and on the common wall its continuous scalar
coefficient is zero.  Applying $D$ and using separate closedness gives
$AB=0$.  Since
\[
 M=v+h_+,\quad m=v+h_-,\quad u=v+h_++h_-,
\]
expansion proves \eqref{r11p:sorting-powers}.  The same expansion,
now using both $AB=0$ and the two scalar-weighted zero identities,
proves \eqref{r11p:sorting-energies}.
\end{proof}

The second identity is useful but has a definite limitation:
polarizing $u(Du)^k$ recovers symmetrized energies, not an arbitrary
ordered scalar-curvature monomial $u_0Du_1\cdots Du_k$.
The parameter argument below supplies precisely that missing step.

\begin{lemma}[Trace of pure PL curvature polynomials]
\label{r11p:pure-trace}
For finite dominant $q:Y\to X$ between normal integral varieties, a
finite sum $C$ of products $Du_1\cdots Du_k$ has an intrinsic actual
trace $\operatorname{Tr}^{\rm curv}_q(C)$, characterized by
\[
 \int_X\operatorname{Tr}^{\rm curv}_q(C)\wedge\eta
                         =\int_YC\wedge q^*\eta
\tag{MT5}\label{r11p:pure-pairing}
\]
for all compact complementary actual tests.  With a parameter, jointly
PL source potentials in fixed source logarithms and $t$ yield a trace
with a product-moment presentation using fixed base logarithms and $t$.
\end{lemma}

\begin{proof}
First let $q$ be Galois.  Sort the orbit of a scalar $u$ by a finite
sorting network.  Repeated use of \eqref{r11p:sorting-powers} gives
\[
                     \sum_{g\in G}(Dg^*u)^k
                            =\sum_{j=1}^d(Du_{(j)})^k.
\]
The ordered potentials $u_{(j)}$ are invariant and descend by
Lemma~\ref{r11p:pl-descent}.  The right side is consequently the
pullback of a downstairs actual form.  Define it to be the trace of
$(Du)^k$.  Polarization,
\[
 Du_1\cdots Du_k
   =\frac1{k!}\sum_{I\subset\{1,\ldots,k\}}
           (-1)^{k-|I|}\left(D\sum_{i\in I}u_i\right)^k,
\]
does the same for all mixed products.  For $k=0$ set
$\operatorname{Tr}^{\rm curv}_q(1)=d$.

This construction satisfies the actual orbit identity
\[
                 q^*\operatorname{Tr}^{\rm curv}_q(C)
                                      =\sum_{g\in G}g^*C.
\]
Wedge with $q^*\eta$, integrate, and use
\eqref{r11p:native-degree} and invariance of native integration under
automorphisms.  Division by $d$ proves \eqref{r11p:pure-pairing}.
Finite pullback faithfulness proves independence of every actual
relation.  Current injectivity gives gluing on ordinary overlaps.
The pairing permits arbitrary actual downstairs multipliers; they
have not been treated as arbitrary source coefficients.

For a general separable finite map take a normal Galois closure
$\widetilde Y\to X$ factoring through $r:\widetilde Y\to Y$, and divide
the Galois trace of $r^*C$ by $e=\deg r$.  Applied to the top form
$Cq^*\eta$, \eqref{r11p:native-degree} proves the required pairing and
independence of the closure.  Characteristic zero covers every map
used here.  All spaces and parameter tori remain over the original
$K$; native generic degrees are not normalized away.

For the last assertion use the parameter version of
Lemma~\ref{r11p:pl-descent} at the sorting step.  The descended
potentials have fixed base moments, so their Hessians and products
have the required product-moment presentation.
\end{proof}

\subsection{Integration along a parameter torus}

\paragraph{The boundary issue.}
Compact analytic support does not imply that an ambient moment
coefficient avoids the tropicalized analytic boundary of an affinoid
chart.  We do not use the single-chart integral formula with that
unverified hypothesis.  Instead we first prove a pushforward identity
for polyhedral currents on the product skeleton, retaining the boundary
exclusion of Proposition 2.38 at every chart calculation.  The
intrinsic skeletal definition of integration then gives the analytic
pairing.  In this subsection $f^\star$ denotes tropical-space pullback
in that proposition, distinguished from the notation $t_c^*$ for the
resulting actual analytic form.

\begin{lemma}[The coordinate product skeleton, including its structures]
\label{r11p:coordinate-product-skeleton}
Let $V$ be a boundaryless ordinary analytic open of pure dimension $n$,
let $b$ be a tuple of analytic units on $V$, and put $S=\Sigma(V,b)$.
The relative Gauss section identifies
\[
 \begin{split}
 \iota_b:S\times\mathbf R&\ \xrightarrow{\ \sim\ }\
       \Sigma(V\times(\mathbf G_m)^{\rm an},(b,z)),\\
                  (y,t)&\longmapsto \nu_t(y),
 \end{split}
\tag{PS1}\label{r11p:product-skeleton-identity}
\]
where $\nu_t(y)$ is the Gauss point of radius $e^{-t}$ in the torus
over $\mathcal H(y)$.  This is an isomorphism of weighted pwl spaces
and of tropical spaces.  Its weight is the product of the native
weight on $S$ and the standard weight on $\mathbf R$, and its linear
function sheaf is generated by the base linear functions and $t$.
The identification is compatible with every ordinary base shrinking
and every additional base moment refinement.
\end{lemma}

\begin{proof}
Let $x$ lie above $y$ in the analytic product.  The graded field
generated by $b$ and $z$ has transcendence degree $n+1$ precisely
when the graded field generated by $b$ at $y$ has transcendence
degree $n$ and $\widetilde z$ is transcendental over
$\widetilde{\mathcal H(y)}$.  Indeed adjoining $z$ increases that
degree by at most one; in the degree-$n$ case
$\widetilde{\mathcal H(y)}$ is algebraic over the graded field
generated by $b$.  This is the graded-transcendence calculation in
(3.27) of the cited proof, here with the single coordinate $z$.

The condition on $\widetilde z$ characterizes the Gauss valuation
over $\mathcal H(y)$: no cancellation among leading monomials is
possible, and for every Laurent polynomial one has
\[
 \left|\sum_i a_i z^i\right|_{\nu_t(y)}
                         =\max_i |a_i|e^{-it}.
\]
Conversely these Gauss valuations have the required graded
transcendence degree.  Thus Definition 3.8 and the absence of analytic
boundary give the claimed bijection.  No valued extension of $K$,
or rationality of the radius in its value group, is needed.

The inverse is $(p,v(z))$ on the displayed skeleton.  It is pwl:
a base skeleton chart is described by an affinoid domain and
valuations of its analytic units; their pullbacks are valuations of
the same units on the product, and $v(z)$ is the extra coordinate.
The corresponding closed chart domains pull back to closed pwl
domains.  This inverse is also proper.  A compact subset of
$S\times\mathbf R$ is contained in a finite union of compact base
analytic domains times a bounded interval.  Its inverse image is
closed in the corresponding finite union of products with compact
analytic annuli, and is therefore compact.  It follows that the
continuous bijective inverse is a homeomorphism.  Subdivision into
local polyhedral charts makes its injective affine pieces invertible
onto their images, so both directions are pwl.  This is the
coordinate specialization of the topology and atlas argument
following (3.31), now without its changes of field or center.

For the weights, take a top-dimensional base cell and a local toric
projection used in Definition 3.8(2) to compute its native degree
$e$.  Its product with the identity on $\mathbf G_m$ is finite flat
of the same degree $e$.  Hence its induced weight is the base weight
times the standard parameter weight.  This calculation retains
residue-field degrees and is exactly the last weight calculation
in the proof of Proposition 3.21.  It is local on cells and hence
determines the entire product weight.

Under the bijection the moments are $(t_b,t)$.
Definition 3.8(3) therefore gives precisely the product linear
function sheaf, not merely a set-theoretic identification.
For additional base moments $b'$, the same Gauss section applies.
The inclusions of base and product skeletons agree under it, and
their weights agree by Remark 3.9.  The new linear functions are the
new base ones together with the same $t$.  This proves the stated
compatibility for all refinements and ordinary restrictions.
\end{proof}

\begin{lemma}[Tropical projection with boundary control]
\label{r11p:tropical-product-push}
Let $S$ be a pure-dimensional tropical space and
$f:S\to\mathbf R^N$ a linear map.  Give $S\times\mathbf R$ the
product weight and linear functions.  Write $p_S$ for its projection
and $\pi:\mathbf R^N\times\mathbf R\to\mathbf R^N$ for the coordinate
projection.  If $\Gamma$ is a Euclidean delta-form with support in
a fixed compact set of parameter values, then
\[
 (p_S)_*\bigl((f,\operatorname{id})^\star\Gamma\bigr)
                                      =f^\star(\pi_*\Gamma)
\tag{PS2}\label{r11p:tropical-product-push-identity}
\]
as polyhedral currents on $S$.  The left side uses only the
proper-on-support pwl-current pushforward of Proposition 2.17.
\end{lemma}

\begin{proof}
The product is a tropical space: on a product polyhedral subdivision,
the vertical balancing conditions are those of $S$, and horizontal
internal facets cancel in pairs with the standard parameter weight.
The support of
$B=(f,\operatorname{id})^\star\Gamma$ is contained in
$S\times J_0$ for a compact interval $J_0$.  Thus it is proper over
$S$, and Proposition 2.17 supplies the polyhedral current
$P=(p_S)_*B$ before any assertion that it is a delta-form.

Work in a polyhedral neighborhood carrying a finite-fiber linear
map $h$ to $\mathbf R^m$, as supplied by the axioms of a sheaf of
linear functions.
Append $f$ and put $g=(f,h)$.  Then $g$ has finite fibers and
$f=A\circ g$, where $A$ is a literal coordinate projection.
Let $H$ be any compact polyhedral subset of the interior of this
neighborhood, and put
\[
 E_H=g(\partial_S H),\qquad
 W_H=\mathbf R^{N+m}\setminus E_H,\qquad
 H_g^\circ=H\setminus g^{-1}(E_H).
\]
The restricted $g:H_g^\circ\to W_H$ is proper.  Write
$C_H=g_*[H_{\dim S},\mu_S]|_{W_H}$.
It is a tropical cycle on $W_H$ by Corollary 2.35.

Choose an interval $J=[a,b]$ containing $J_0$ in its interior.
The boundary of $H\times J$ is contained in
\[
        (\partial_S H\times J)\ \cup\ (H\times\{a,b\}).
\]
Apply the characterizing identity (2.40) of Proposition 2.38 to
this compact product and the refinement $g\times\operatorname{id}$.
On the boundary-free region $W_H\times(a,b)$ it gives
\[
 (g\times\operatorname{id})_*
       \bigl(B|_{H_g^\circ\times(a,b)}\bigr)
       =(C_H\times\mathbf R)\wedge
                             (A\times\operatorname{id})^*\Gamma.
\tag{PS3}\label{r11p:boundary-chart-product}
\]
Both sides vanish near the artificial parameter endpoints because
$\Gamma$ does.  They may therefore be viewed on $W_H\times\mathbf R$,
but the exclusion of $E_H$ in the base remains in force.  There
is no assertion of this chart identity across $g(\partial_S H)$.

Let $\pi_g:W_H\times\mathbf R\to W_H$ be projection.  Functoriality
of pwl-current pushforward and the Euclidean projection formula
give, entirely on $W_H$,
\[
\begin{split}
  g_*\bigl(P|_{H_g^\circ}\bigr)
   &=\pi_{g,*}\!\left((C_H\times\mathbf R)\wedge
                          (A\times\operatorname{id})^*\Gamma\right)\\
   &=C_H\wedge
                  \pi_{g,*}(A\times\operatorname{id})^*\Gamma\\
   &=C_H\wedge A^*(\pi_*\Gamma).
\end{split}
\tag{PS4}\label{r11p:boundary-chart-push}
\]
The last step needs only coordinate Fubini.  For a coordinate
projection $A$, its Euclidean delta-pullback adds free standard
coordinates.  On each weighted carrier of $\Gamma$, fiber
integration in $t$ before or after adding those coordinates
has the same fibers, coefficient integral and quotient weights.
Thus
$\pi_{g,*}(A\times\operatorname{id})^*\Gamma=A^*(\pi_*\Gamma)$.
This carrierwise argument applies to an arbitrary Euclidean
delta-form, with compact parameter support; it does not invoke a
general affine base-change lemma.

The last expression in \eqref{r11p:boundary-chart-push} is exactly
the characterizing expression (2.40) for
$f^\star(\pi_*\Gamma)$.  Varying the compact polyhedral set $H$
inside the fixed finite-fiber $g$-chart, Lemma 2.37 shows that
$P=f^\star(\pi_*\Gamma)$ on that chart.  These charts cover $S$,
which proves \eqref{r11p:tropical-product-push-identity}.
In particular the polyhedral pushforward has the asserted
delta-form representative; this is a conclusion of the calculation.
\end{proof}

\begin{lemma}[Integration on the product-moment subspace]
\label{r11p:product-integration}
Let $p:U\times(\mathbf G_m)^{\rm an}\to U$ be projection.  Suppose
\[
                   \beta=t_{(c,z)}^*\Gamma
\]
for a fixed base analytic-unit tuple $c$, $t=v(z)$, and a Euclidean
delta-form $\Gamma(u,t)$ supported in a fixed compact set of $t$.
With $\pi(u,t)=u$, define
\[
                       I_p(\beta)=t_c^*(\pi_*\Gamma).
\tag{MT6}\label{r11p:product-definition}
\]
This is intrinsic among product-moment presentations and satisfies
\[
             \int_UI_p(\beta)\wedge\eta
                    =\int_{U\times\mathbf G_m}\beta\wedge p^*\eta
\tag{MT7}\label{r11p:product-pairing}
\]
for every compact actual complementary test.  It commutes with
$d',d''$ on these presentations.
\end{lemma}

\begin{proof}
\textit{1. A skeletal identity before integration.}
Euclidean pushforward supplies the actual candidate
$\theta=t_c^*(\pi_*\Gamma)$.
On an arbitrary ordinary $V\subseteq U$, refine the base moments
by tuple concatenation, $b=(c,h)$, and put $S=\Sigma(V,b)$.
Under \eqref{r11p:product-skeleton-identity}, the restriction of
$\beta$ to the product skeleton is
$(t_c|_S,\operatorname{id})^\star\Gamma$.
Lemma~\ref{r11p:tropical-product-push} therefore gives
\[
       \theta|_S
        =(p_S)_*\bigl(\beta|_{S\times\mathbf R}\bigr).
\tag{PS5}\label{r11p:actual-product-skeletal}
\]
This is an equality of intrinsic polyhedral currents on $S$ for
every base refinement.  Its proof used the boundary-excluded chart
identities \eqref{r11p:boundary-chart-product}--\eqref{r11p:boundary-chart-push},
not an integral over a global tropicalization of $V$.

If two product-moment presentations define the same actual $\beta$,
concatenate both base tuples and any further test tuple.  Their
restrictions to every resulting product skeleton agree by Definition
4.2.  Equation \eqref{r11p:actual-product-skeletal} makes their
candidate outputs agree on every corresponding base skeleton.
Lemma 4.4, using all refinements, gives actual equality.  This proves
independence of presentations and their relations, and also their
compatibility on ordinary overlaps, before any current pairing.

\textit{2. Products with arbitrary actual base forms.}
On an open $V$ where an actual form $\eta$ has a moment presentation,
take $b$ to be the concatenation of its moments and $c$, and write
$\eta=t_b^*\eta_b$.  Let $A$ be the coordinate projection selecting
$c$.  The product $\beta\wedge p^*\eta$ has Euclidean coefficient
\[
           (A\times\operatorname{id})^*\Gamma\wedge\pi_b^*\eta_b
\]
in the fixed moments $(b,z)$, with the same compact parameter support.
Euclidean projection and coordinate Fubini give
\[
 \pi_{b,*}\!\left(
      (A\times\operatorname{id})^*\Gamma\wedge\pi_b^*\eta_b\right)
                         = A^*(\pi_*\Gamma)\wedge\eta_b.
\]
Apply \eqref{r11p:actual-product-skeletal} to this coefficient.
For $S=\Sigma(V,b)$ it gives
\[
 (p_S)_*
     \left((\beta\wedge p^*\eta)|_{S\times\mathbf R}\right)
                         =(\theta\wedge\eta)|_S.
\tag{PS6}\label{r11p:product-test-skeletal}
\]
The coefficient $\eta_b$ here is an arbitrary Euclidean delta-form.
No restriction to smooth tests or to products of model curvatures
has occurred.

\textit{3. Intrinsic skeletal integration and ordinary partitions.}
Now let $\eta$ be compact and of complementary bidegree.  Cover its
support by finitely many ordinary opens $V_i$ on which it has the
presentations used in the preceding step, and choose an
original-smooth partition $(\lambda_i)$ subordinate to this cover
on a neighborhood of the support.  Put
$S_i=\Sigma(V_i,b_i)$.
On $V_i$ the unweighted form $\theta\wedge\eta$ is top-degree
and is presented in $b_i$; on $V_i\times\mathbf G_m$ the
unweighted form $\beta\wedge p^*\eta$ is top-degree and is
presented in $(b_i,z)$.  Corollary 4.9(2) puts their supports in
$S_i$ and $S_i\times\mathbf R$, respectively.

Multiplication by $\lambda_i$ or $p^*\lambda_i$ cannot enlarge
these supports.  Their paired supports are compact: the relevant
base support is compact inside $V_i$, and the parameter values lie
in a fixed compact interval, hence in a compact analytic annulus.
Finitely many affinoid neighborhoods of the base support give
compact total paired support.  This compactness justifies the
integrals; it is not used to exclude any tropicalized base boundary.

The restriction $\lambda_i|_{S_i}$ need only be PWS.  Remark 4.3(1)
makes its multiplication coefficientwise, and Proposition 2.17
gives the coefficientwise projection formula
\[
 (p_{S_i})_*\!\left(p_{S_i}^*(\lambda_i|_{S_i})\,T\right)
                  =(\lambda_i|_{S_i})(p_{S_i})_*T
\]
for the relevant polyhedral current $T$.
Use Definition 4.10(1) to evaluate the analytic top integrals on
these containing skeletons, followed by (2.17) and the integrable
current pushforward formula in Definition 4.10(4).  With
\eqref{r11p:product-test-skeletal}, this gives
\[
\begin{split}
 \int_{U\times\mathbf G_m}p^*\lambda_i\,\beta\wedge p^*\eta
 &=\int_{S_i\times\mathbf R}
      p_{S_i}^*(\lambda_i|_{S_i})
          (\beta\wedge p^*\eta)|_{S_i\times\mathbf R}\\
 &=\int_{S_i}(\lambda_i|_{S_i})(\theta\wedge\eta)|_{S_i}\\
 &=\int_U\lambda_i\,\theta\wedge\eta.
\end{split}
\]
Definition 4.10(2) permits summing this identity over the ordinary
partition.  It proves \eqref{r11p:product-pairing} for every
compact actual test.  The partition is not required to have a
single Euclidean presentation in the preassigned tuple $b_i$.

Finally, Euclidean proper-on-support pushforward commutes with
$d',d''$, and differentiation retains the compact parameter support.
Applying this to \eqref{r11p:product-definition}, whose actual
independence is already proved, gives the asserted differential
compatibility.  Equivalently it follows from the linear-map
current identity of Definition 2.29(3) in
\eqref{r11p:actual-product-skeletal}.
\end{proof}

\begin{lemma}[One-parameter evaluation of a nonlinear PL scalar]
\label{r11p:one-parameter}
Let $u$ be a bounded PL scalar on an ordinary open $V$.  Choose
$\chi\in C_c^\infty(\mathbf R)$ equal to one near its range.  On
$V\times(\mathbf G_m)^{\rm an}$ put
\[
                         H=\max(u-t,0).
\]
For every actual form $C$ on $V$ with a common moment presentation,
\[
                  I_p\bigl(\chi(t)t\,DH\wedge p^*C\bigr)=uC.
\tag{MT8}\label{r11p:graph-evaluation}
\]
Only one parameter occurs, independently of the number of source
moments needed to present $u$.
\end{lemma}

\begin{proof}
Write $u=\varphi(\ell)$ for a finite Euclidean PL scalar.  Work first
on its Euclidean moment domain with $H(x,t)=\max(\varphi(x)-t,0)$.
On each cell where $\varphi$ is affine, the graph face $t=\varphi(x)$
in $DH$ has projection weight one.  The coefficient of $t$ in the
defining affine difference is $-1$, so projection of its calibrated
divisor measure is the standard base measure.  Weighting by $\chi(t)t$
gives $\chi(\varphi(x))\varphi(x)$.

There are also vertical faces above walls of $\varphi$, in the region
$t<\varphi(x)$.  It would be false to identify $DH$ with the graph
alone.  On every such vertical face, however, projection drops the
dimension of its pure carrier; its pure Euclidean pushforward is
zero, exactly as in Mihatsch's Example 2.5.  Thus
\[
               \pi_*\bigl(\chi(t)t\,D\max(\varphi-t,0)\bigr)
                                  =\chi(\varphi)\varphi.
\]
The affine-cell equality extends across walls as an equality of
actual scalar functions; the right side is continuous and the
Euclidean pushforward is a delta-form of bidegree zero.

The affine projection formula now permits an arbitrary delta-form
$C$:
\[
 \pi_*\bigl(\chi(t)t\,DH\wedge\pi^*C\bigr)
       =\pi_*(\chi(t)t\,DH)\wedge C.
\]
Apply Lemma~\ref{r11p:product-integration} after moment refinement.
Only after actual moment pullback do we use the cutoff choice:
$\chi(\varphi(\ell))\varphi(\ell)=\chi(u)u=u$.
The ambient extension $\varphi$ need not be bounded, so
$\chi(\varphi)\varphi=\varphi$ is not asserted on its entire
Euclidean domain.
This argument has not pulled $C$ through the nonlinear PL graph;
only the affine projection pullback and projection formula are used.
\end{proof}

\subsection{The minimal finite trace theorem}

\begin{theorem}[Finite trace on the scalar-curvature module]
\label{r11p:minimal-trace-theorem}
Let $q:Y\to X$ be finite dominant between normal integral separated
varieties over $K$.  Every $\omega\in\mathscr M^k(Y^{\rm an})$ has a
unique actual trace on $X^{\rm an}$ satisfying
\[
                 \int_X\operatorname{Tr}^{\mathscr M}_q(\omega)
                                      \wedge\eta
                          =\int_Y\omega\wedge q^*\eta
\tag{MT9}\label{r11p:minimal-pairing}
\]
for all compact complementary actual tests.  It is local on the
ordinary topology, respects actual relations, and retains the full
native generic degree.  For a locally presented generator
$\omega=u_0Du_1\cdots Du_k$, a local formula is
\[
 \boxed{\displaystyle
 \operatorname{Tr}^{\mathscr M}_q(\omega)
 = I_{p_X}\!\left(
     \chi(t)t\,
     \operatorname{Tr}^{\rm curv}_{q\times\operatorname{id}}
       \left(D\max(u_0-t,0)\wedge\prod_{i=1}^kDu_i\right)
                         \right).}
\tag{MT10}\label{r11p:trace-formula}
\]
\end{theorem}

\begin{proof}
Shrink around an entire finite source fiber so that finitely many
generator presentations are combined as in
Lemma~\ref{r11p:pl-descent}, and $u_0$ is bounded there.  Choose
$\chi=1$ near that whole range.  The curvature trace inside
\eqref{r11p:trace-formula} exists by Lemma~\ref{r11p:pure-trace}.
Its parameter-uniform assertion supplies a product-moment presentation,
so multiplication by $\chi(t)t$ and application of $I_{p_X}$ are
legitimate.

Pair with an arbitrary compact actual test $\eta$.  First
Lemma~\ref{r11p:product-integration}, then
\eqref{r11p:pure-pairing} with the actual downstairs multiplier
$\chi(t)t\,p_X^*\eta$, and finally
Lemma~\ref{r11p:one-parameter}, give
\[
\begin{split}
 \int_X\operatorname{Tr}^{\mathscr M}_q(\omega)\eta
 &=\int_{X\times\mathbf G_m}\chi(t)t\,
     \operatorname{Tr}^{\rm curv}_{q\times\operatorname{id}}
        (DH\prod_iDu_i)\,p_X^*\eta\\
 &=\int_{Y\times\mathbf G_m}\chi(t)t\,DH\prod_iDu_i\,
                                          p_Y^*q^*\eta\\
 &=\int_Yu_0\prod_iDu_i\,q^*\eta.
\end{split}
\]
All these integrals have compact paired support.  The identity is
proved for all actual tests, not merely for model or original-smooth
tests.  Add generators.  If an expression represents zero as an
actual source form, its trace has zero pairing with every test, and
is zero by Proposition 4.11.  The same argument removes dependence
on the cutoff, source moments, quotient moments, Galois closure, and
all local choices.  It also glues the local answers in the ordinary
topology and proves uniqueness.
\end{proof}

